% =====================================================================
%  Resumen ejecutivo / Abstract
% =====================================================================
\section*{Resumen}
\addcontentsline{toc}{section}{Resumen}

El material de apoyo de un curso universitario se acumula a lo largo de varios
semestres hasta formar corpus en los que conviven versiones redundantes,
afirmaciones contradictorias y contenidos desactualizados. Mantener la coherencia
de ese corpus es una tarea de curación que recae sobre el docente y que la
literatura sobre carga académica sitúa entre las actividades poco visibles en los
modelos institucionales de asignación de trabajo \parencite{kenny2023life}. La
incorporación de inteligencia artificial generativa no resuelve el problema por sí
sola: la evidencia sistemática disponible indica que la IA generativa
\emph{redistribuye} la carga docente más de lo que la reduce, porque añade trabajo
de revisión, verificación y juicio sobre sus propias salidas
\parencite{su2026efficiency}.

Este documento presenta \textbf{EduCurator}, un sistema de curación asistida de
documentación de cursos que combina recuperación de evidencia sobre el corpus del
propio curso, detectores de redundancia e inconsistencia, un subagente con
herramientas, evidencia estructurada y revisión humana obligatoria antes de
cualquier publicación. La investigación se formula bajo
\emph{Design Science Research} \parencite{peffers2007design,hevner2004design}, con
punto de entrada en la actividad de diseño y desarrollo: el artefacto ya existe
como MVP funcional y el trabajo pendiente es su demostración y evaluación
controlada.

El documento cumple dos funciones simultáneas. Como \textbf{protocolo}, define la
pregunta de investigación, las proposiciones, las variables operacionalizadas, el
benchmark con defectos plantados y controles negativos, la batería de
configuraciones de comparación (heurísticas, recuperación dispersa, recuperación
densa, canalización híbrida y sistema completo), las métricas de detección,
\emph{grounding}, calibración y carga de trabajo, y el plan de análisis
estadístico. Como \textbf{reporte de avance}, delimita qué está construido, qué
evidencia técnica existe hoy y qué permanece abierto.

\begin{estado}
\small
\textbf{Alcance de lo que aquí se afirma.} Este documento \emph{no} reporta
resultados experimentales. No se afirma que EduCurator detecte inconsistencias con
una precisión determinada, que reduzca el tiempo de curación, que sus puntajes de
confianza estén calibrados ni que los docentes lo perciban como útil. Esas
afirmaciones requieren la evidencia que la Parte V está diseñada para producir y
que todavía no se ha recogido.
\end{estado}

\vspace{0.4em}
\noindent\textbf{Palabras clave:} curación de contenidos educativos; generación
aumentada por recuperación; agentes basados en modelos de lenguaje;
\emph{human-in-the-loop}; trazabilidad; \emph{design science research}; carga de
trabajo docente.

\vspace{1.2em}

\section*{Abstract}
\addcontentsline{toc}{section}{Abstract}

Course support material accumulates over successive semesters into corpora where
redundant versions, contradictory statements and outdated content coexist. Keeping
such a corpus coherent is a curation task that falls on the instructor and that the
academic-workload literature places among the activities institutional workload
models tend to overlook \parencite{kenny2023life}. Introducing generative AI does
not solve this by itself: systematic evidence indicates that generative AI
\emph{redistributes} teaching workload more than it reduces it, since it adds
review, verification and judgement work over its own outputs
\parencite{su2026efficiency}.

This document presents \textbf{EduCurator}, an assisted curation system for course
documentation that combines evidence retrieval over the course corpus, redundancy
and inconsistency detectors, a tool-using sub-agent, structured evidence, and
mandatory human review before any publication. The research is framed as
\emph{Design Science Research} \parencite{peffers2007design,hevner2004design} with a
design- and development-centred entry point: the artefact already exists as a
functional MVP, and what remains is its demonstration and controlled evaluation.

The document works both as a \textbf{protocol} --- research question, propositions,
operationalised variables, a benchmark with planted defects and negative controls,
a ladder of comparison configurations, detection, grounding, calibration and
workload metrics, and a statistical analysis plan --- and as a
\textbf{progress report} that states what has been built, what technical evidence
exists today, and what remains open. \textbf{No experimental results are reported
here.}

\vspace{0.4em}
\noindent\textbf{Keywords:} educational content curation; retrieval-augmented
generation; LLM agents; human-in-the-loop; traceability; design science research;
academic workload.
